% Exported from Math++. Compile with pdflatex.
\documentclass[11pt,a4paper]{article}
\usepackage[T1]{fontenc}
\usepackage[utf8]{inputenc}
\usepackage{lmodern}
\usepackage{amsmath,amssymb,amsthm,mathtools}
\usepackage{graphicx,booktabs,xcolor}
\usepackage[hidelinks]{hyperref}
\usepackage[margin=25mm]{geometry}
\providecommand{\degree}{\ensuremath{^\circ}}
\providecommand{\xmark}{\ensuremath{\times}}
\providecommand{\omicron}{o}
\providecommand{\oiint}{\mathop{\oint\mkern-12mu\oint}}
\theoremstyle{plain}
\newtheorem{theorem}{Theorem}
\newtheorem{lemma}[theorem]{Lemma}
\newtheorem{corollary}[theorem]{Corollary}
\newtheorem{proposition}[theorem]{Proposition}
\newtheorem{conjecture}[theorem]{Conjecture}
\theoremstyle{definition}
\newtheorem{definition}[theorem]{Definition}
\newtheorem{example}[theorem]{Example}
\theoremstyle{remark}
\newtheorem*{remark}{Remark}
\setlength{\parindent}{0pt}\setlength{\parskip}{0.6em}
\title{Assumptions}
\date{}
\begin{document}
\maketitle
Without knowing anything about $x$, Math++ keeps the safe form and says why in the working (a query that ends in \texttt{?!} shows it):

\begin{equation*}
  \sqrt{x^2} = \left|x\right|
\end{equation*}
A power integral has two cases, so the answer is piecewise:

\begin{align*}
  \int x^n \, dx = \begin{cases} \frac{x^{n + 1}}{n + 1} & n \neq -1 \\ \ln\left(\left|x\right|\right) & \text{otherwise} \end{cases} + C
\end{align*}
An improper integral that depends on a parameter:

\begin{align*}
  \int_1^\infty x^{-p} \, dx = \begin{cases} \infty & p \le 1 \\ -\frac{1}{-p + 1} & \text{otherwise} \end{cases}
\end{align*}
A condition in brackets holds for its block only:

\begin{equation*}
  \sqrt{y^2} = -y \quad (y < 0)
\end{equation*}
From here on, $x$ is positive, $n$ is an integer and $a$ is negative.

Assume $x > 0$.

Assume $n \in \mathbb{Z}$, $a < 0$.

\begin{equation*}
  \sqrt{x^2} = x
\end{equation*}
\begin{equation*}
  \ln(e^x) = x
\end{equation*}
\begin{equation*}
  \cos(n \pi) = \left(-1\right)^{n}
\end{equation*}
\begin{equation*}
  \lim_{x \to \infty} e^{a x} = 0
\end{equation*}
\begin{equation*}
  \int_0^\infty e^{a x} \, dx = -\frac{1}{a}
\end{equation*}
Assume $p > 1$.

\begin{equation*}
  \int_1^\infty x^{-p} \, dx = -\frac{1}{-p + 1}
\end{equation*}
\section*{Mistakes the checker catches}

Assumptions that contradict the ones above are an error:

Assume $x < 0$.

{\small\color{red!60!black}$\times$ these assumptions can't all hold: x > 0 and x < 0}

\end{document}
